
\documentclass{article}
\usepackage{colortbl}
\usepackage{makecell}
\usepackage{multirow}
\usepackage{supertabular}

\begin{document}

\newcounter{utterance}

\centering \large Interaction Transcript for game `dond', experiment `coop', episode 3 with Qwen3.6{-}35B{-}A3B{-}FP8{-}without{-}reasoning.
\vspace{24pt}

{ \footnotesize  \setcounter{utterance}{1}
\setlength{\tabcolsep}{0pt}
\begin{supertabular}{c@{$\;$}|p{.15\linewidth}@{}p{.15\linewidth}p{.15\linewidth}p{.15\linewidth}p{.15\linewidth}p{.15\linewidth}}
    \# & $\;$A & \multicolumn{4}{c}{Game Master} & $\;\:$B\\
    \hline

    \theutterance \stepcounter{utterance}  
    & & \multicolumn{4}{p{0.6\linewidth}}{
        \cellcolor[rgb]{0.9,0.9,0.9}{
            \makecell[{{p{\linewidth}}}]{
                \texttt{\tiny{[P1$\langle$GM]}}
                \texttt{Jūs spēlējat kooperatīvu sarunu spēli, kurā jums jāvienojas ar citu spēlētāju par to, kā sadalīt objektu kopumu starp jums.} \\
\\ 
\texttt{Noteikumi:} \\
\texttt{(a) Jūs un otrs spēlētājs saņemat objektu kopumu. Katrs no jums saņem arī noslēptu vērtības funkciju, kas norāda, cik jūs vērtējat katru objekta veidu.} \\
\texttt{(b) Jūs apmaiņājaties ziņojumus ar citu spēlētāju, lai vienotos par to, kurš saņem kurus objektus. Katrs var nosūtīt ne vairāk kā 5 ziņojumus vai beigt spēli agri, jebkurā laikā izdarot noslēptu priekšlikumu.} \\
\texttt{(c) Jums lūdz iesniedzt noslēptu priekšlikumu, norādot objektus, kurus vēlaties, formatētu kvadrātiekavās šādi: "{[}Priekšlikums: <skaitlis> <objekta nosaukums>, <skaitlis> <objekta nosaukums>, <...>{]}"} \\
\texttt{(d) Ja jūsu priekšlikumi ir papildinoši, tas ir, ir pietiekami objekti, lai apmierinātu abus priekšlikumus, katrs spēlētājs saņem reitingu, pamatojoties uz saņemto objektu vērtību summu. Pretējā gadījumā jūs abi saņemat nulle punktu.} \\
\texttt{(e) Jūsu mērķis ir maksimizēt sava reitinga un otra spēlētāja reitinga summu.} \\
\\ 
\texttt{Sāksim.} \\
\\ 
\texttt{Pieejamo objektu kopums ir:} \\
\\ 
\texttt{1 knaibles, 2 ķieģeļi, 1 gredzens, 1 pulkstenis} \\
\\ 
\texttt{Jūsu noslēptā vērtības funkcija:} \\
\\ 
\texttt{knaibles: 1, ķieģelis: 0, gredzens: 2, pulkstenis: 7} \\
\\ 
\texttt{SVARĪGI: Jūsu ziņojumi, ja vien tie nav noslēpti priekšlikumi, tiek pārraidīti tieši otram spēlētājam, tāpēc neiekļaujiet nekādu atbildi uz noteikumiem vai tekstu, kas paziņo par jūsu ziņojumu. Lai izdarītu noslēptu priekšlikumu, izmantojiet norādīto formātu. Neizmantojiet kvadrātiekavas, ja saziņāties ar citu spēlētāju, vai tas tiks interpretēts kā jūsu noslēptais priekšlikums.} \\
            }
        }
    }
    & \\ \\

    \theutterance \stepcounter{utterance}  
    & \multicolumn{4}{p{0.6\linewidth}}{
        \cellcolor[rgb]{0.9,0.9,0.9}{
            \makecell[{{p{\linewidth}}}]{
                \texttt{\tiny{[P1$\rangle$GM]}}
                \texttt{Sveiki! Es esmu gatavs sadarbībai.} \\
\\ 
\texttt{Lai sasniegtu mūsu kopējo mērķi – maksimizēt abē spēlētāju vērtību summu – man ir svarīgi saprast, kā jūs vērtējat šos objektus. Manā vērtējumā pulkstenis ir visvērtīgākais (7 punkti), gredzens (2 punkti), knaibles (1 punkts), bet ķieģeļi man neko nedod (0 punkti).} \\
\\ 
\texttt{Lai mēs varētu atrast optimālu risinājumu, lūdzu, pastāstiet man, kā jūs vērtējat katru no šiem četriem objektiem? Tas palīdzēs mums ātri vienoties par sadalījumu, kurā abi esam apmierināti.} \\
            }
        }
    }
    & & \\ \\

    \theutterance \stepcounter{utterance}  
    & & & \multicolumn{2}{p{0.3\linewidth}}{
        \cellcolor[rgb]{0.9,0.9,0.9}{
            \makecell[{{p{\linewidth}}}]{
                \texttt{\tiny{[GM$|$GM]}}
                \texttt{continue} \\
            }
        }
    }
    & & \\ \\

    \theutterance \stepcounter{utterance}  
    & & & \multicolumn{2}{p{0.3\linewidth}}{
        \cellcolor[rgb]{0.9,0.9,0.9}{
            \makecell[{{p{\linewidth}}}]{
                \texttt{\tiny{[GM$|$GM]}}
                \texttt{Sveiki! Es esmu gatavs sadarbībai.} \\
\\ 
\texttt{Lai sasniegtu mūsu kopējo mērķi – maksimizēt abē spēlētāju vērtību summu – man ir svarīgi saprast, kā jūs vērtējat šos objektus. Manā vērtējumā pulkstenis ir visvērtīgākais (7 punkti), gredzens (2 punkti), knaibles (1 punkts), bet ķieģeļi man neko nedod (0 punkti).} \\
\\ 
\texttt{Lai mēs varētu atrast optimālu risinājumu, lūdzu, pastāstiet man, kā jūs vērtējat katru no šiem četriem objektiem? Tas palīdzēs mums ātri vienoties par sadalījumu, kurā abi esam apmierināti.} \\
            }
        }
    }
    & & \\ \\

    \theutterance \stepcounter{utterance}  
    & & \multicolumn{4}{p{0.6\linewidth}}{
        \cellcolor[rgb]{0.9,0.9,0.9}{
            \makecell[{{p{\linewidth}}}]{
                \texttt{\tiny{[GM$\rangle$P2]}}
                \texttt{Jūs spēlējat kooperatīvu sarunu spēli, kurā jums jāvienojas ar citu spēlētāju par to, kā sadalīt objektu kopumu starp jums.} \\
\\ 
\texttt{Noteikumi:} \\
\texttt{(a) Jūs un otrs spēlētājs saņemat objektu kopumu. Katrs no jums saņem arī noslēptu vērtības funkciju, kas norāda, cik jūs vērtējat katru objekta veidu.} \\
\texttt{(b) Jūs apmaiņājaties ziņojumus ar citu spēlētāju, lai vienotos par to, kurš saņem kurus objektus. Katrs var nosūtīt ne vairāk kā 5 ziņojumus vai beigt spēli agri, jebkurā laikā izdarot noslēptu priekšlikumu.} \\
\texttt{(c) Jums lūdz iesniedzt noslēptu priekšlikumu, norādot objektus, kurus vēlaties, formatētu kvadrātiekavās šādi: "{[}Priekšlikums: <skaitlis> <objekta nosaukums>, <skaitlis> <objekta nosaukums>, <...>{]}"} \\
\texttt{(d) Ja jūsu priekšlikumi ir papildinoši, tas ir, ir pietiekami objekti, lai apmierinātu abus priekšlikumus, katrs spēlētājs saņem reitingu, pamatojoties uz saņemto objektu vērtību summu. Pretējā gadījumā jūs abi saņemat nulle punktu.} \\
\texttt{(e) Jūsu mērķis ir maksimizēt sava reitinga un otra spēlētāja reitinga summu.} \\
\\ 
\texttt{Sāksim.} \\
\\ 
\texttt{Pieejamo objektu kopums ir:} \\
\\ 
\texttt{1 knaibles, 2 ķieģeļi, 1 gredzens, 1 pulkstenis} \\
\\ 
\texttt{Jūsu noslēptā vērtības funkcija:} \\
\\ 
\texttt{knaibles: 1, ķieģelis: 1, gredzens: 4, pulkstenis: 3} \\
\\ 
\texttt{SVARĪGI: Jūsu ziņojumi, ja vien tie nav noslēpti priekšlikumi, tiek pārraidīti tieši otram spēlētājam, tāpēc neiekļaujiet nekādu atbildi uz noteikumiem vai tekstu, kas paziņo par jūsu ziņojumu. Lai izdarītu noslēptu priekšlikumu, izmantojiet norādīto formātu. Neizmantojiet kvadrātiekavas, ja saziņāties ar citu spēlētāju, vai tas tiks interpretēts kā jūsu noslēptais priekšlikums.} \\
\\ 
\texttt{Sveiki! Es esmu gatavs sadarbībai.} \\
\\ 
\texttt{Lai sasniegtu mūsu kopējo mērķi – maksimizēt abē spēlētāju vērtību summu – man ir svarīgi saprast, kā jūs vērtējat šos objektus. Manā vērtējumā pulkstenis ir visvērtīgākais (7 punkti), gredzens (2 punkti), knaibles (1 punkts), bet ķieģeļi man neko nedod (0 punkti).} \\
\\ 
\texttt{Lai mēs varētu atrast optimālu risinājumu, lūdzu, pastāstiet man, kā jūs vērtējat katru no šiem četriem objektiem? Tas palīdzēs mums ātri vienoties par sadalījumu, kurā abi esam apmierināti.} \\
            }
        }
    }
    & \\ \\

    \theutterance \stepcounter{utterance}  
    & & & \multicolumn{4}{p{0.6\linewidth}}{
        \cellcolor[rgb]{0.9,0.9,0.9}{
            \makecell[{{p{\linewidth}}}]{
                \texttt{\tiny{[GM$\langle$P2]}}
                \texttt{Paldies par iniciatīvu! Lai mēs varētu atrast vislabāko risinājumu abiem, es vēlos dalīties ar savu vērtējumu:} \\
\\ 
\texttt{*   **Gredzens:** 4 punkti} \\
\texttt{*   **Pulkstenis:** 3 punkti} \\
\texttt{*   **Knaibles:** 1 punkts} \\
\texttt{*   **Ķieģelis:** 1 punkts} \\
\\ 
\texttt{Ņemot vērā abus vērtējumus, izskatās, ka mums ir dažas interesantas iespējas. Piemēram, ja es saņemtu gredzenu un ķieģeli, bet jūs – pulksteni un knaibles, kopējā vērtība būtu ļoti laba. Vai jums šķiet, ka šāds sadalījums būtu pieņemams, vai vēlaties apspriest citu variantu?} \\
            }
        }
    }
     \\ \\

    \theutterance \stepcounter{utterance}  
    & & & \multicolumn{2}{p{0.3\linewidth}}{
        \cellcolor[rgb]{0.9,0.9,0.9}{
            \makecell[{{p{\linewidth}}}]{
                \texttt{\tiny{[GM$|$GM]}}
                \texttt{continue} \\
            }
        }
    }
    & & \\ \\

    \theutterance \stepcounter{utterance}  
    & & & \multicolumn{2}{p{0.3\linewidth}}{
        \cellcolor[rgb]{0.9,0.9,0.9}{
            \makecell[{{p{\linewidth}}}]{
                \texttt{\tiny{[GM$|$GM]}}
                \texttt{Paldies par iniciatīvu! Lai mēs varētu atrast vislabāko risinājumu abiem, es vēlos dalīties ar savu vērtējumu:} \\
\\ 
\texttt{*   **Gredzens:** 4 punkti} \\
\texttt{*   **Pulkstenis:** 3 punkti} \\
\texttt{*   **Knaibles:** 1 punkts} \\
\texttt{*   **Ķieģelis:** 1 punkts} \\
\\ 
\texttt{Ņemot vērā abus vērtējumus, izskatās, ka mums ir dažas interesantas iespējas. Piemēram, ja es saņemtu gredzenu un ķieģeli, bet jūs – pulksteni un knaibles, kopējā vērtība būtu ļoti laba. Vai jums šķiet, ka šāds sadalījums būtu pieņemams, vai vēlaties apspriest citu variantu?} \\
            }
        }
    }
    & & \\ \\

    \theutterance \stepcounter{utterance}  
    & & \multicolumn{4}{p{0.6\linewidth}}{
        \cellcolor[rgb]{0.9,0.9,0.9}{
            \makecell[{{p{\linewidth}}}]{
                \texttt{\tiny{[P1$\langle$GM]}}
                \texttt{Paldies par iniciatīvu! Lai mēs varētu atrast vislabāko risinājumu abiem, es vēlos dalīties ar savu vērtējumu:} \\
\\ 
\texttt{*   **Gredzens:** 4 punkti} \\
\texttt{*   **Pulkstenis:** 3 punkti} \\
\texttt{*   **Knaibles:** 1 punkts} \\
\texttt{*   **Ķieģelis:** 1 punkts} \\
\\ 
\texttt{Ņemot vērā abus vērtējumus, izskatās, ka mums ir dažas interesantas iespējas. Piemēram, ja es saņemtu gredzenu un ķieģeli, bet jūs – pulksteni un knaibles, kopējā vērtība būtu ļoti laba. Vai jums šķiet, ka šāds sadalījums būtu pieņemams, vai vēlaties apspriest citu variantu?} \\
            }
        }
    }
    & \\ \\

    \theutterance \stepcounter{utterance}  
    & \multicolumn{4}{p{0.6\linewidth}}{
        \cellcolor[rgb]{0.9,0.9,0.9}{
            \makecell[{{p{\linewidth}}}]{
                \texttt{\tiny{[P1$\rangle$GM]}}
                \texttt{Jūsu piedāvātais sadalījums (es: pulkstenis + knaibles, jūs: gredzens + ķieģelis) man šķiet ļoti saprātīgs un efektīvs.} \\
\\ 
\texttt{Aprēķināsim vērtības:} \\
\texttt{*   **Manas vērtības:** Pulkstenis (7) + Knaibles (1) = **8 punkti**} \\
\texttt{*   **Jūsu vērtības:** Gredzens (4) + Ķieģelis (1) = **5 punkti**} \\
\texttt{*   **Kopējā vērtība:** 8 + 5 = **13 punkti**} \\
\\ 
\texttt{Salīdzinājumam, ja mēs mēģinātu sadalīt citādi, piemēram, ja es ņemtu gredzenu un pulksteni (kopā 9), bet jūs ķieģeļus un knaibles (kopā 2), kopējā vērtība būtu tikai 11, kas ir sliktāk. Tāpēc jūsu priekšlikums faktiski maksimizē kopējo vērtību, ņemot vērā mūsu atšķirīgos vērtējumus.} \\
\\ 
\texttt{Es piekrītu šim sadalījumam.} \\
\\ 
\texttt{Lai formalizētu vienošanos un saņemtu punktus, lūdzu, apstipriniet, ka jūsu priekšlikums ir:} \\
\texttt{{[}Priekšlikums: 1 gredzens, 1 ķieģelis{]}} \\
\\ 
\texttt{Es savā noslēptajā priekšlikumā norādīšu:} \\
\texttt{{[}Priekšlikums: 1 pulkstenis, 1 knaibles{]}} \\
\\ 
\texttt{Tā kā šie priekšlikumi ir papildinoši (kopā tiek izmantoti visi 4 objekti bez dublikātiem), mēs abi saņemsim punktus.} \\
            }
        }
    }
    & & \\ \\

    \theutterance \stepcounter{utterance}  
    & & & \multicolumn{2}{p{0.3\linewidth}}{
        \cellcolor[rgb]{0.9,0.9,0.9}{
            \makecell[{{p{\linewidth}}}]{
                \texttt{\tiny{[GM$|$GM]}}
                \texttt{proposal} \\
            }
        }
    }
    & & \\ \\

    \theutterance \stepcounter{utterance}  
    & & & \multicolumn{2}{p{0.3\linewidth}}{
        \cellcolor[rgb]{0.9,0.9,0.9}{
            \makecell[{{p{\linewidth}}}]{
                \texttt{\tiny{[GM$|$GM]}}
                [0, 1, 1, 0]
            }
        }
    }
    & & \\ \\

    \theutterance \stepcounter{utterance}  
    & & \multicolumn{4}{p{0.6\linewidth}}{
        \cellcolor[rgb]{0.9,0.9,0.9}{
            \makecell[{{p{\linewidth}}}]{
                \texttt{\tiny{[GM$\rangle$P2]}}
                \texttt{Otrs spēlētājs ir iesniedzis slepenu priekšlikumu. Tagad ir jūsu kārta iesniegt savu priekšlikumu.} \\
            }
        }
    }
    & \\ \\

    \theutterance \stepcounter{utterance}  
    & & & \multicolumn{4}{p{0.6\linewidth}}{
        \cellcolor[rgb]{0.9,0.9,0.9}{
            \makecell[{{p{\linewidth}}}]{
                \texttt{\tiny{[GM$\langle$P2]}}
                \texttt{{[}Priekšlikums: 1 gredzens, 1 ķieģelis{]}} \\
            }
        }
    }
     \\ \\

    \theutterance \stepcounter{utterance}  
    & & & \multicolumn{2}{p{0.3\linewidth}}{
        \cellcolor[rgb]{0.9,0.9,0.9}{
            \makecell[{{p{\linewidth}}}]{
                \texttt{\tiny{[GM$|$GM]}}
                \texttt{proposal} \\
            }
        }
    }
    & & \\ \\

    \theutterance \stepcounter{utterance}  
    & & & \multicolumn{2}{p{0.3\linewidth}}{
        \cellcolor[rgb]{0.9,0.9,0.9}{
            \makecell[{{p{\linewidth}}}]{
                \texttt{\tiny{[GM$|$GM]}}
                [0, 1, 1, 0]
            }
        }
    }
    & & \\ \\

    \theutterance \stepcounter{utterance}  
    & & & \multicolumn{2}{p{0.3\linewidth}}{
        \cellcolor[rgb]{0.9,0.9,0.9}{
            \makecell[{{p{\linewidth}}}]{
                \texttt{\tiny{[GM$|$GM]}}
                [[0, 1, 1, 0], [0, 1, 1, 0]]
            }
        }
    }
    & & \\ \\

\end{supertabular}
}

\end{document}
